\section{Where the pattern and the condition come from}
\label{sec:ddl}

\subsection{The question the comprehension cannot answer}

The comprehension of \S\ref{sec:cfl} presupposes two things it does not supply.
Its pattern \texttt{@spec} must have a grammar of terms to match against. Its condition \texttt{Chi(spec)} is a claim about
structure, so there must be a notion of what structure a specimen has.

In SQL both come from the schema. \texttt{spec} is a row of a known table with
known column types, and \texttt{chi} is a predicate over those columns. The
schema is the data definition language, and the query language leans on it
entirely.

Rholang's schema is the \texttt{language!} block, and the reason to keep reading
is that it is a great deal more than a schema.

\subsection{A \texttt{language!} block, whole}

Here is one, complete~\cite{mettailrust}. It defines the \rhoc{}
calculus~\cite{MR05} itself --- which is a
useful thing to see early, because it shows that the calculus rholang is built
on is not privileged: it is one language definition among others, expressible in
the same block a developer would use for an invoice format.

\begin{lstlisting}[language=mettail]
language! {
    name: RhoCalc,

    types {
        Proc
        Name
    },

    terms {
        PZero  . |- "0" : Proc;
        PDrop  . n:Name |- "*" "(" n ")" : Proc;
        POutput. n:Name, q:Proc |- n "!" "(" q ")" : Proc;
        PInput . n:Name, ^x.p:[Name -> Proc] |- n "?" x "." "{" p "}" : Proc;
        PPar   . ps:HashBag(Proc) |- "{" ps.*sep("|") "}" : Proc;
        NQuote . p:Proc |- "@" "(" p ")" : Name;
    },

    equations {
        (NQuote (PDrop N)) = N;
    },

    rewrites {
        (PPar {(PInput n ^x.p), (POutput n q), ...rest})
            ~> (PPar {(subst ^x.p (NQuote q)), ...rest});

        (PDrop (NQuote P)) ~> P;

        if S ~> T then (PPar {S, ...rest}) ~> (PPar {T, ...rest});
    },
}
\end{lstlisting}

Four blocks. \texttt{types} names the syntactic categories. \texttt{terms} gives
one constructor per line, each with its arguments, its concrete syntax, and its
result type --- so the grammar and the pretty-printer come from the same
declaration. \texttt{equations} says which distinct terms are to be regarded as
the same. \texttt{rewrites} says how terms move --- in the style Plotkin introduced for
operational semantics~\cite{Plotkin81} and Milner used to such effect for
processes~\cite{Milner92}.

From that block, the toolchain generates the abstract syntax types, a parser and
a bidirectional printer, capture-avoiding substitution that respects binding,
and a rewrite engine.

\subsection{The three rungs}

The four blocks form a ladder, and where a given language definition stops on
that ladder determines what kind of thing it is.

\begin{center}
\small
\begin{tabular}{@{}L{0.26\textwidth}L{0.32\textwidth}L{0.32\textwidth}@{}}
\toprule
\textbf{Blocks populated} & \textbf{What you can present} & \textbf{The familiar name} \\
\midrule
\texttt{types}, \texttt{terms} &
free term algebras &
algebraic data types \\[3pt]
$+$ \texttt{equations} &
quotients of free algebras &
finitely presentable universal algebra \\[3pt]
$+$ \texttt{rewrites} &
theories with dynamics &
domain-specific languages \\
\bottomrule
\end{tabular}
\end{center}

An ordinary data definition language stops on the first rung. Records, sums,
products, lists, maps: that is the whole of what most type systems can say about
data, and it is a lot. The second and third rungs are where this differs, and
the second rung is the one that will surprise a working developer most, so take
it first.

\subsection{Rung two: equations, or the comment your type system cannot read}

Every codebase of any size contains comments of this form:

\begin{quote}\small
\emph{This is a list, but order doesn't matter --- always compare as a set.}\\
\emph{This is a string, but compare case-insensitively.}\\
\emph{Addresses with an empty line-2 are equal to the same address without it.}
\end{quote}

Each is an equation. Each is being maintained by hand, in prose, because the
type system cannot express it. The consequences are familiar: a hand-written
\texttt{equals}, a normalisation function called on the way in and forgotten in
one code path, a cache keyed by the wrong representative, and a class of bug
that only appears when two spellings of the same value meet.

An \texttt{equations} block moves that comment into the definition:

\begin{lstlisting}[language=mettail]
equations {
    (NQuote (PDrop N)) = N;
}
\end{lstlisting}

The payoff is not that the compiler now checks something. It is that
\emph{pattern matching happens up to the equations}. When the comprehension of
\S\ref{sec:cfl} matches \texttt{@spec} against a datum, it matches modulo the
theory that types the channel. A developer never writes normalisation, never
chooses a canonical form, and never has the bug where one code path forgot.

This is why the \texttt{PPar} constructor above takes a \texttt{HashBag} rather
than a list: parallel composition is associative and commutative, and rather
than being a special case hard-wired into an interpreter, that is a property of
the data declaration. The same mechanism is available for an invoice.

\begin{aside}[The reach of the second rung]
``Finitely presentable universal algebra'' is the mathematician's name for the
second rung, and it is worth translating: it means anything you can specify by
saying what the constructors are and which combinations of them are equal.
Monoids, groups, lattices, sets-as-quotients-of-lists, normalised paths,
case-insensitive strings, and commutative parallel composition are all on this
rung. Most of a data model's invisible invariants live here.
\end{aside}

\subsection{Rung three: rewrites, or data that wiggles}

The third rung adds \texttt{rewrites}, and with it the thing an ordinary schema
cannot do: the data moves on its own.

\begin{lstlisting}[language=mettail]
rewrites {
    Beta      . |- (App (Lam fun) arg) ~> (eval fun arg);
    AppCongL  . | M0 ~> M1 |- (App M0 N) ~> (App M1 N);
}
\end{lstlisting}

That is the lambda calculus, in two of the four lines a real definition takes.
The first line is the computation rule; the second is a congruence rule saying
where the first is allowed to fire. Together they make a term of type
\texttt{Term} not an inert value but a small running program.

Once a channel is typed by such a definition, sending on that channel means
sending a term of a language --- a term that can be inspected, matched against,
and stepped by the recipient. Two agents that share a language definition can
exchange programs in a domain-specific language constructed for the purpose, and
the recipient knows what it received because the definition came with it.

The slogan for this rung is \emph{the data can wiggle}, and its practical form
is that the boundary between data and code, which most systems maintain by
convention and defend with a serialisation format, is a dial here rather than a
wall.

\subsection{Channels are typed by language definitions}

The connection back to \S\ref{sec:cfl} is one sentence: \emph{a channel's type
is a language definition}.

That single design decision explains the rest of the note. Sending on the
channel means sending a term of that theory. Matching a pattern on it means
matching a term of that theory, up to that theory's equations. And --- the point
of \S\ref{sec:logic} --- asking a condition about what arrived means asking a
question in a logic that is \emph{generated from} that theory, because the
theory says what there is to ask about.

For the forager, this means the specimen language is a language definition,
\texttt{Chi} is a formula in the logic that definition generates, and there was
never a moment at which someone hand-designed a predicate vocabulary.

\subsection{MeTTa, and what this adds to it}
\label{sec:metta}

Readers coming from SingularityNET's Hyperon stack~\cite{Hyperon} will have
recognised the shape of the third rung, because MeTTa is built on it. The relationship is worth
stating precisely, in both directions.

\begin{center}
\small
\begin{tabular}{@{}L{0.24\textwidth}L{0.30\textwidth}L{0.38\textwidth}@{}}
\toprule
\textbf{MeTTa} & \textbf{MeTTaIL} & \textbf{Relationship} \\
\midrule
rules &
the \texttt{rewrites} block &
subsumed; MeTTa's rules are rewrites in this sense \\[3pt]
atom spaces &
channels &
subsumed, and typed: a channel carries terms of a declared language, an
atom space carries atoms \\[3pt]
--- &
the \texttt{equations} block &
\textbf{absent in MeTTa}: no equational layer, so every quotient must be
hand-coded as rules, and confluence becomes the programmer's problem \\[3pt]
query, then update &
one witnessed communication &
\textbf{absent in MeTTa}: query and update are separate acts, so there is no
atomic multi-party meeting and no witness \\
\bottomrule
\end{tabular}
\end{center}

The two gaps compound. Without equations, a developer who wants ``these two
atoms are the same'' writes rewrite rules that say so, and then owns the
confluence question --- whether the order in which rules fire can change the
answer. Without a witnessed transaction, a developer who wants two atom spaces
updated together writes a protocol, and then owns the partial-failure question.
Both are exactly the questions the two blocks and the communication rule
answer for free.

None of which is a dismissal. MeTTa has a large working ecosystem, a substantial
body of practice around the atom space as a shared knowledge medium, and
Probabilistic Logic Networks~\cite{PLN}, which is a serious piece of engineering for
reasoning under uncertainty. It is worth saying plainly that PLN is not left
behind here: in \S\ref{sec:graded} it comes back as one of the value algebras a
\texttt{where} clause can be graded in, which is to say that PLN inference can
drive execution rather than sitting beside it. The relationship this note is
proposing is embedding with an addition, not replacement.

\subsection{What an operation costs}
\label{sec:cost}

One more piece belongs to the data definition story, because it is the thing
that makes the forager's budget real rather than notional.

Cost accounting in this system adds exactly two surface forms, and both of them
lower to ordinary rholang:

\begin{center}
\small
\begin{tabular}{@{}lll@{}}
\toprule
\textbf{Surface} & \textbf{Lowers to} & \textbf{Which is} \\
\midrule
\texttt{s :: ()} & \texttt{$\Sigma\sem{s}$!(Nil)} & a send: mint a token \\
\texttt{\{\% P \%\}[s]} & \texttt{for(t <- $\Sigma\sem{s}$)\{*t | P\}} & a receipt: run \texttt{P} once a token arrives \\
\bottomrule
\end{tabular}
\end{center}

That is the entire mechanic. Spending fuel \emph{is} a communication: the gate
consumes one token, releases its payload, and then \texttt{P} runs. If no token
is available on the signature channel, the gate parks forever and \texttt{P}
never runs --- and that is what metering means here. A budget is a stack of
nested sends, and its depth is how many operations it funds.

Two things follow that a developer should notice.

First, metering is not a virtual-machine feature bolted onto a language. It is a
program \emph{in} the language, and the transpiler's job is only to write it for
you. Anything true of ordinary rholang is true of metered rholang.

Second, it composes with Claim~\ref{claim:witness} in a satisfying way. A
communication that is not funded is a communication that does not happen, and it
is indistinguishable from a communication whose guard failed. Both are
non-events. There is no separate failure mode for running out of budget, no
exception, and nothing to unwind --- which is exactly the property that lets the
forager's death be modelled as ``it stops'' rather than as an error path.

The repository~\cite{f1r3node} ships a worked demonstration of this at scale: a supply chain in
which factories produce, carriers ship, sellers wholesale and buyers retail,
with money and inventory both conserved, and every single operation gated by a
token. It is the best available answer to ``show me this doing something real''.

\subsection{What actually runs}
\label{sec:status}

Honesty about status is more useful to a developer than enthusiasm, so:

\begin{center}
\small
\begin{tabular}{@{}L{0.40\textwidth}L{0.52\textwidth}@{}}
\toprule
\textbf{Capability} & \textbf{Status} \\
\midrule
\texttt{language!} blocks: types, terms, equations, rewrites; generated parser,
printer, substitution, rewrite engine &
implemented in \texttt{mettail-rust}~\cite{mettailrust}, with language
definitions shipped for
lambda, ambient, the \rhoc{} calculus, a space calculus and a calculator \\[3pt]
\texttt{where} guards on receipts, and \texttt{where} guards on \texttt{match}
cases with fall-through &
implemented on the \texttt{cost-accounting-transpiler} branch of
\texttt{f1r3node-rust}~\cite{f1r3node}, with runnable examples and compile tests \\[3pt]
Cost accounting: minting, gates, compound signatures, budget stacks &
implemented on the same branch, with a transpiler and a documented demo \\[3pt]
Conditions drawn from the \emph{generated logic} (\S\ref{sec:logic}) rather than
from ordinary boolean expressions &
partial. The guard slot exists and is evaluated before commit; the structural
and behavioural connectives are specified in the research notes and are being
brought into the branch \\[3pt]
Graded conditions and simulation (\S\ref{sec:graded}) &
implemented as a separate simulator crate~\cite{WeightedGSLT}, with the
stochastic and quantum readings both built and tested; not part of the node's execution path, by
design (\S\ref{sec:simnotinterp}) \\
\bottomrule
\end{tabular}
\end{center}

The condition language a developer can write \emph{today} is boolean expressions
over the variables the patterns bound --- \texttt{x > 0}, and so on. Everything
in the next section is about what that slot is going to hold, and why the choice
is more interesting than it looks.
